% !TeX root = ../papers/nc-p3.tex
\section{正次元ファイバーを許す四移動と反復可能な分岐の排除}
\label{sec:recurrent-point-correspondences}
位置 $i$ を固定し、標準四移動で $V_{i+4}=V_i$ とする。
点は $\mathbb P(V_i^*)$ の任意の幾何点でよい。
$p,q$ に対し
\[
 d(p,q)=\dim\frac{A_{i+1,i+4}}
 {R_i(p) A_{i+2,i+4}+A_{i+1,i+3}R_{i+3}(q)}
\]
を定める。前者の評価は関係の先頭、後者は末尾で取る。
$A_{i+1,i+4}$ は有限Frobenius候補で次元20である。この余核の正次元条件は
両点の対の上で二群の列がそれぞれ $p,q$ に線型な $20\times120$ 行列の階数条件である。
その被約閉対応を $D$ とする。実際の四辺の点列が $p,q$ を結べば必ず $(p,q)\in D$ である。
命題\ref{prop:cut-arrow-growth}は $d(p,p)$ を使う場合に当たる。

\begin{theorem}\label{thm:recurrent-point-correspondences}
標準四窓では以下の三条件が全ての位置で成り立つ。
\begin{enumerate}
\item 任意の有限個の相異なる点を頂点に、$p$ から $q$ に $d(p,q)$ 本の有向辺を置く。
各強連結成分のうち巡回路を持つものは、全ての辺の多重度が1の単純有向巡回路である。
一つの有向道に沿って互いに異なるそのような成分を通過できる個数は4以下である。
点を一般位置に置く仮定はない。
\item 閉既約部分多様体の周期列 $Y_0,\ldots,Y_l=Y_0$ と
$D\cap(Y_j\times Y_{j+1})$ 内の被約既約対応 $B_j$ があり、
両射影がそれぞれの $Y$ に支配的なら、各 $B_j$ の両射影は双有理である。
また $d=1$ が各 $B_j$ の一般点で成立する。
特にこのような周期列では正次元の一般ファイバーも二枚以上の分離被覆も生じない。
\item 同じ既約 $Y$ から $Y$ に戻る二つのこのような対応列があるとする。
長さを反復で揃えると、その二列は一般の始点からの各中間点まで一致する。
従って同じ $Y$ 上の相異なる支配的自己対応を二つ同時に使う反復分岐もない。
\end{enumerate}
\end{theorem}
\begin{proof}
上の余核の双対空間を $K(p,q)$ と書く。三文字テンソルがこの空間に属する条件は、
両内側の関係の零化と、その左右に $p,q$ を付けた二つの境界関係の零化である。
従って $p_0,\ldots,p_m$ の辺列について
\[
 p_0\otimes K(p_0,p_1)\otimes p_1\otimes\cdots
        \otimes K(p_{m-1},p_m)\otimes p_m
 \ \subset\ (A_{i,i+4m+1})^*
\]
となる。点の代表の倍率は独立性に影響しない。
一つの始点から長さ $m$ の二分木を作り、分かれる二辺の終点が違うか、
同じ終点でも $K$ の選んだテンソルが線型独立なら、その全ての葉の積テンソルは独立である。
実際、各辺を三文字テンソルと終点の計四文字で読むと、最初の分岐の二つの四文字は独立である。
それぞれの片方だけを消す線型形式で縮約し、残りの木に帰納する。
各頂点で一分岐だけの部分と、長さが等しい二本の異なる道を一単位として読む場合にも適用できる。
後者では中間点が違えばその点を置いた因子で区別できる。
全頂点を同時に一次独立に取る必要がない。

有限強連結成分が単純な巡回路でなければ、ある頂点に異なる二つの出辺がある。
それぞれに同じ頂点への帰路を付け、反復して正の長さを揃える。
二つの道の四文字ブロックの積テンソルは最初の出辺で区別できる。
任意の順に繰り返せば長さに対して指数個の独立な元ができ、三次のHilbert増大に反する。
次に有向道が $s$ 個の巡回成分を通るとする。相異なる成分間の固定した連結路を選び、
各巡回成分では退出頂点でその巡回路を回ってから次へ出る。
各巡回長の公倍数を $b$ とし、巡回部にさらに $b n_j$ 辺ずつを挿入する。
$\sum n_j=N$ で得た $\binom{N+s-1}{s-1}$ 個の道は同じ長さであり、二分岐の木をなす。
分かれる最初の頂点で退出辺と留まる辺が異なる。従って対応する積テンソルは独立であり、
$N$ に関する三次増大から $s\le4$ である。

支配的対応列では、所定の次側稠密開集合に入る点を各段で選べる。
もし一つの出射影の一般ファイバーが正次元または幾何的に二点以上なら、
どの有限反復回数 $m$ についても、一般始点からその段で毎周二点に分かれる木が作れる。
詳しくは帰納的に終点側の稠密開を引き戻し、対応の稠密開の像で始点側の稠密開を縮める。
選択肢が二つ必要な段では対応の自己ファイバー積から対角を除く。
支配性と一般ファイバーの条件によって、その二つの点をともに指定の稠密開に取れる。
各 $m$ で始点が変わっても、得られる $2^m$ 個の独立な元は同じ次数の経路空間に入る。
従って一般ファイバーは一点であり、標数零より双有理である。
逆向きにも同じ成長制限を使える。
一般の辺で $d\ge2$ の場合は中間点を動かす代わりに $K$ の二独立元を選べば同じ矛盾となる。

最後に二つの列を同じ長さにしたとき、初めて異なる中間点への二つの双有理写像は
稠密開集合上で異なる値を与える。両方の終点の写像はなお支配的なので、
一方か他方を反復ごとに選ぶ上記の二分木が構成できる。
\end{proof}

\paragraph{有限障害としての利用。}
例えば $p\ne q$ で $d(p,p),d(p,q),d(q,p)>0$ が同時に成り立つ点対は、
二つの戻り道による除外を与える。任意の固定した有限有向多重グラフを指定しても、
全頂点の相異条件と必要な次数の小行列式だけでその出現を射影消去で判定できる。
また相異なる点が一次独立な範囲なら、木を二分岐に限らず各 $K$ の基底を同時に使え、
$M=(d(p_a,p_b))$ として $\dim A_{i,i+4m+1}\ge\sum_{a,b}(M^m)_{ab}$ となる。
成長次数だけに基づく巡回成分数の上限4は改善できない。
四文字の非可換単項式代数で降順の隣接二文字を零とすれば、基底は弱増加語である。
座標四点の $M_{ab}$ は $a\le b$ で $\binom{b-a+3}{3}$、ほかは零となり、
四個のループを順に通過でき、上の下界は実際の三次Hilbert関数に一致する。
この単項式代数を標準AS例とする主張は用いていない。
